% !PRINTED 291
\weberpage{291}
%<<CONT
also ist die Charakteristik gleich dem halben Ueberschuss der ersten Zahl über die zweite.

Diesem Satz können wir folgenden Ausdruck geben, wobei er von Grössenverhältnissen gänzlich unabhängig erscheint und nur von den Lagenverhältnissen der Curven und ihrer Schnittpunkte abhängt.

Unterscheiden wir die Schnittpunkte von $\varphi$, $\psi$, je nachdem sie Austrittspunkte $A(\varphi,\psi)$ oder Eintrittspunkte $E(\varphi,\psi)$ sind, durch das Zeichen $\alpha$ und $\varepsilon$, ebenso, je nachdem sie äussere oder innere Punkte zu der Curve $f$ sind, durch $\alpha$ und $\varepsilon$, geben dann den Punkten $\alpha,\varepsilon$ und $\varepsilon,\alpha$ den \gesperrt{Charakter $+1$, den Punkten $\alpha,\alpha$ und $\varepsilon,\varepsilon$ den Charakter $-1$}\textsuperscript{1)}, \gesperrt{so ist die Charakteristik des Functionensystems $(f,\varphi,\psi)$ gleich der halben Summe der Charaktere der sämmtlichen Schnittpunkte von $\varphi,\psi$.}

Die Sätze und Begriffsbestimmungen, die wir hier gegeben haben, bleiben unverändert bestehen, auch wenn die betrachteten Curven aus mehreren in sich geschlossenen Zügen bestehen; auch können Doppelpunkte bei den Curven vorhanden sein; es muss nur von jeder Fortschrittsrichtung auf einer der Curven völlig bestimmt sein, ob sie als positiv oder als negativ zu betrachten ist, d.\,h. es muss jedes Stück einer Curve $f$ oder $\varphi$ oder $\psi$ Flächentheile von einander trennen, in denen die Functionen $f$ oder $\varphi$ oder $\psi$ entgegengesetzte Vorzeichen haben.

Auszuschliessen sind nur die Fälle, in denen eine der Curven durch einen Doppelpunkt der anderen hindurchgeht, oder in denen die Curven einander berühren, oder in denen die drei Curven $f$, $\varphi$, $\psi$ durch \gesperrt{einen} Punkt gehen; denn in diesen Fällen würde die Bestimmung eines Punktes als Eintrittspunkt oder Austrittspunkt zweifelhaft werden.

Erstrecken sich die Curven ins Unendliche, so muss man sie, um unsere Sätze anwendbar zu machen, durch willkürlich hinzugefügte Curvenstücke abschliessen, wie wir im nächsten Paragraphen sehen werden.

\noindent\rule{4cm}{0.4pt}

\noindent\textsuperscript{1)} \gesperrt{Kronecker} versteht unter Charakter eines Punktes etwas Anderes.
